\chapter{Limitations and Future Work}
\label{ch:limitations}

\section{Methodological and Empirical Limitations}
While this project provides a rigorous, controlled investigation into reinforcement learning for portfolio management, several realistic limitations must be acknowledged:

\subsection{Passive Benchmark Outperformance and Survivorship Bias}
The equal-weight 1/N strategy demonstrated superior cumulative profit and Sharpe ratio compared to all deep RL agents. This is attributable to two primary factors:
\begin{enumerate}
    \item \textbf{Survivorship Bias in Equity Selection:} The 10 stocks selected (e.g., Reliance, TCS, HDFC Bank, ICICI Bank) represent established corporate leaders that survived and prospered over the 2011--2025 window. In an asset universe characterised by strong upward drift, a passive, fully-invested allocation captures market beta at negligible transaction costs ($0.61\%$).
    \item \textbf{Cash Drag from Learned Risk Aversion:} The continuous Dirichlet policy learned to maintain cash buffers (up to $32.4\%$ in M4) to guard against volatility. While this successfully curbed maximum drawdown to $-16.04\%$, it penalized nominal capital compounding during bull markets.
\end{enumerate}

\subsection{Market Regime Dependence and Single-Epoch Hyperparameter Selection}
As disclosed in Chapter~\ref{ch:results}, the turnover regularisation coefficient $\lambda_{\text{turn}} = 0.1$ was tuned strictly on Fold 1 validation across two seeds. Subsequent folds encompassed vastly different market regimes—notably the COVID-19 crash of March 2020 (Fold 2) and the global inflationary cycle of 2022 (Fold 3). A constant hyperparameter tuned in a calm market proved overly restrictive during volatile recovery phases where rapid rebalancing was necessary.

\subsection{Simplified Execution and Market Impact Assumptions}
The environment models execution costs as a fixed proportional fee $c_{\text{trans}} = 0.001$ ($10\text{ bps}$). While appropriate for retail and liquid institutional execution, this model neglects:
\begin{itemize}
    \item Non-linear market impact and price slippage for large order sizes.
    \item Intraday execution dynamics (trades are executed at the official daily closing price).
    \item Short-selling constraints and margin financing costs.
\end{itemize}

\subsection{Computational and Stochastic Architecture Constraints}
Due to computational constraints, policies were restricted to a single-layer LSTM with 512 hidden units, trained for 300,000 steps per run. Furthermore, continuous Dirichlet sampling exhibits inherent stochasticity, leading to noticeable variance across random seeds (e.g., Fold 2 return std $\pm 14.8\%$ in M4).

\section{Directions for Future Research}
The insights gained from this study motivate several promising avenues for future inquiry:
\begin{enumerate}
    \item \textbf{Dynamic and Adaptive Turnover Regularisation:} Rather than maintaining a static penalty $\lambda_{\text{turn}}$, future models could employ adaptive penalty scheduling conditioned on real-time market volatility (e.g., VIX) or learned meta-policy networks.
    \item \textbf{Attention and Transformer Backbones:} Replacing recurrent LSTM cells with temporal multi-head attention mechanisms (such as Temporal Fusion Transformers) could enhance the extraction of long-range cross-asset correlations.
    \item \textbf{Macroeconomic and Alternative Feature Integration:} Augmenting technical indicators with macroeconomic variables (interest rates, crude oil prices, USD/INR foreign exchange rates) and sentiment features from financial disclosures.
    \item \textbf{Dynamic Cash Hurdles:} Incorporating an active risk-free interest rate yield on the cash reserve to mitigate the harsh cash drag observed in equity bull markets.
\end{enumerate}
